\documentclass[11pt,a4paper]{article}
\usepackage[T1]{fontenc}
\usepackage[utf8]{inputenc}
\usepackage{amsmath,amsthm,mathtools}
\usepackage{newpxtext,newpxmath}
\usepackage[margin=27mm,headheight=15pt]{geometry}
\usepackage{microtype}
\usepackage{booktabs,array,longtable}
\usepackage{xcolor}
\usepackage{enumitem}
\usepackage{fancyhdr}
\usepackage{aliascnt}
\usepackage{titlesec}
\usepackage{listings}
\usepackage{hyperref}
\usepackage[nameinlink,noabbrev]{cleveref}
\definecolor{Olive}{HTML}{596446}
\definecolor{Ink}{HTML}{253029}
\definecolor{Pale}{HTML}{F2F4EC}
\definecolor{Muted}{HTML}{62685E}
\hypersetup{colorlinks=true,linkcolor=Olive,citecolor=Olive,urlcolor=Olive,
 pdftitle={Friendly Order Types Through Incomparability Components},
 pdfauthor={Research manuscript prepared with ChatGPT},
 pdfsubject={Finite friendly ranks, successor defects, and disjoint sums}}
\titleformat{\section}{\large\bfseries\color{Olive}}{\thesection}{0.65em}{}
\titleformat{\subsection}{\normalsize\bfseries\color{Ink}}{\thesubsection}{0.6em}{}
\pagestyle{fancy}
\fancyhf{}
\fancyhead[L]{\small\color{Muted}Friendly order types}
\fancyhead[R]{\small\color{Muted}Incomparability components and finite tails}
\fancyfoot[C]{\small\thepage}
\renewcommand{\headrulewidth}{0.3pt}
\setlength{\parskip}{3pt}
\setlength{\parindent}{1.15em}
\setlist{itemsep=3pt,topsep=5pt}
\emergencystretch=1.5em
\newtheorem{theorem}{Theorem}[section]
\newaliascnt{lemma}{theorem}
\newtheorem{lemma}[lemma]{Lemma}
\aliascntresetthe{lemma}
\crefname{lemma}{lemma}{lemmas}
\newaliascnt{proposition}{theorem}
\newtheorem{proposition}[proposition]{Proposition}
\aliascntresetthe{proposition}
\crefname{proposition}{proposition}{propositions}
\newaliascnt{corollary}{theorem}
\newtheorem{corollary}[corollary]{Corollary}
\aliascntresetthe{corollary}
\crefname{corollary}{corollary}{corollaries}
\theoremstyle{definition}
\newaliascnt{definition}{theorem}
\newtheorem{definition}[definition]{Definition}
\aliascntresetthe{definition}
\crefname{definition}{definition}{definitions}
\newaliascnt{example}{theorem}
\newtheorem{example}[example]{Example}
\aliascntresetthe{example}
\crefname{example}{example}{examples}
\theoremstyle{remark}
\newaliascnt{remark}{theorem}
\newtheorem{remark}[remark]{Remark}
\aliascntresetthe{remark}
\crefname{remark}{remark}{remarks}
\DeclareMathOperator{\mot}{o}
\DeclareMathOperator{\fot}{f}
\DeclareMathOperator{\str}{str}
\DeclareMathOperator{\Bad}{Bad}
\DeclareMathOperator{\rk}{rk}
\DeclareMathOperator{\Max}{Max}
\newcommand{\core}[1]{{#1}^{\circ}}
\newcommand{\tail}[1]{T(#1)}
\newcommand{\inc}{\mathbin{\perp}}
\newcommand{\ds}{\mathbin{\sqcup}}
\newcommand{\Mr}{\mathcal M^{r}}
\newcommand{\N}{\mathbb N}
\newcommand{\eps}{\varepsilon}
\newcommand{\smallstatus}[1]{\par\noindent\colorbox{Pale}{\parbox{\dimexpr\linewidth-2\fboxsep\relax}{\small #1}}\par}
\lstset{basicstyle=\small\ttfamily,breaklines=true,frame=single,
 rulecolor=\color{Olive},backgroundcolor=\color{Pale},columns=fullflexible,
 keepspaces=true,showstringspaces=false}

\begin{document}
\hypersetup{pageanchor=false}
\begin{titlepage}
\vspace*{13mm}
{\color{Olive}\large ORDER THEORY / RESEARCH MANUSCRIPT\par}
\vspace{9mm}
{\huge\bfseries\color{Ink}Friendly Order Types Through\par
Incomparability Components\par}
\vspace{5mm}
{\Large Finite ranks, successor defects,\par and corrected disjoint-sum laws\par}
\vspace{10mm}
{\large Prepared for Vladimir Reshetnikov\par}
{\normalsize 19 September 2026\par}
\vspace{10mm}
\begin{abstract}
The friendly order type of a well-partial order is the rank of its tree of bad
sequences whose successive choices retain an incomparable witness. We attack
the problem of computing this invariant structurally and compositionally.
For every finite poset $P$, we prove
\[
\fot(P)=|P|-c(P),
\]
where $c(P)$ is the number of connected components of its incomparability
graph. The proof gives an optimal witness sequence and a polynomial-time
algorithm. For an infinite incomparability-connected well-partial order, we
prove that the friendly order type is either its maximal order type or, when
that type is a successor, its immediate predecessor. A canonical finite
upper tail and a saturation test on the remaining core decide between the
two cases. This yields exact disjoint-sum laws, including equality with
maximal order type whenever both summands are infinite.

We also give a direct counterexample to a disjoint-sum identity printed in a
2023 paper. More strongly, two explicit countable orders have identical
maximal order type, height, width, and friendly order type, but adjoining an
incomparable singleton gives different friendly order types. Thus those four
invariants do not support a general disjoint-sum calculation. The archive
includes independent finite-rank and graph implementations, checked on all
101,660 naturally labeled posets with at most seven vertices.
\end{abstract}
\vfill
\smallstatus{\textbf{Status.} This is a research manuscript with detailed proofs,
not a peer-reviewed or proof-assistant-verified publication. The finite and
explicit counterexample arguments are elementary. The general transfinite
classification uses an identified, previously published limit-floor bound.
No claim of independently established priority is made. The broader problem
of effective computation for arbitrary presented well-partial orders is not
claimed to be settled.}
\end{titlepage}
\hypersetup{pageanchor=true}
\pagenumbering{roman}
\tableofcontents
\newpage
\pagenumbering{arabic}

\section{The problem and the outcome}\label{sec:problem}

Vialard's 2024 thesis asks: ``Can we compute the fot compositionally?''\
\cite[p.~109]{Vialard2024}. Here \emph{fot} means
\emph{friendly order type}. The invariant had appeared in the author's MFCS
2023 work, where it determines the ordinal width of the multiset ordering
\cite[Theorem~3.4]{Vialard2023}.

We extract a definite, tractable part of that research direction: find an
intrinsic finite formula, determine the obstruction at infinite successor
ranks, and understand disjoint sums. We also ask whether the standard three
ordinal invariants, augmented by the friendly order type itself, contain
enough information for such a calculation. This last question has a negative
answer, witnessed by concrete countable orders below. The structural question
has a positive answer in terms of incomparability components, maximal order
types, and a single bit attached to an infinite core.

The most useful conclusions are the following. First, a finite friendly rank
is simply the number of edges in a spanning forest of the incomparability
graph. Second, an infinite connected component has a defect of at most one
from maximal order type; the defect is completely classified. Third,
disjoint sums of two infinite orders have no defect at all. Fourth, adding
one isolated-in-the-order point can distinguish inputs that agree in all
four familiar ordinal invariants.

These assertions are not all new definitions of special-purpose ranks: the
rank being calculated is exactly Vialard's friendly order type. We write
$\fot(P)$ instead of $\mot_{\perp}(P)$ to keep the formulas readable.

\subsection{What is and is not an open-status claim}

The documented source poses a broad research question, not the exact finite
equality proved here as a separately numbered conjecture. We do not relabel
our extracted subproblem as a previously published conjecture. The current
check covered the author's thesis and publication page, the published MFCS
paper, and targeted searches for later work on the friendly order type. No
matching component-and-core formula was located. Search coverage and the
absence of a located result do not establish priority.

There is a relevant version distinction. Proposition~4.1(2) of the 2023
published paper prints an unrestricted disjoint-sum identity. The 2024 thesis
instead gives a sufficient limit-ordinal condition in Theorem~5.4.8(2).
Section~\ref{sec:printed} supplies a direct counterexample to the earlier
printed identity; it is not presented as contradicting the later conditional
statement. The rest of the manuscript provides replacement formulas and
makes clear which prior results are used.

\section{Definitions, ranks, and arithmetic}\label{sec:definitions}

All orders are sets, and the ambient foundational setting is ZFC. A
\emph{well-partial order}, or wpo, is a partial order having no infinite bad
sequence. A finite sequence $x_0,\ldots,x_{k-1}$ is \emph{bad} when
$x_i\not\leq x_j$ for every $i<j$. For $x\in P$, put
\[
\uparrow_P x=\{y\in P:x\leq y\},\qquad
P_x=P\setminus\uparrow_P x.
\]
Thus selecting $x$ in a bad sequence removes $x$ and all its successors,
not its predecessors. For a finite history $s$, the residual is
$P_s=P\setminus\bigcup_{x\in s}\uparrow_Px$.

The rank of a node in a well-founded tree is the supremum of the ranks of
its immediate successors plus one. A leaf has rank zero. The maximal order
type $\mot(P)$ is the rank of the bad-sequence tree, equivalently the maximum
order type of a linear extension. The existence of a maximizing linear
extension and the standard finite-sum laws are established wpo theory
\cite{deJonghParikh1977,DSS2020}. In particular,
\begin{equation}\label{eq:mot-rec}
\mot(P)=\sup_{x\in P}\bigl(\mot(P_x)+1\bigr),\qquad
\mot(A\ds B)=\mot(A)\oplus\mot(B).
\end{equation}
Here $A\ds B$ is a disjoint sum: every element of $A$ is incomparable to
every element of $B$. By contrast, $A+B$ is an ordinal sum, with every
element of $A$ below every element of $B$.

The symbol $+$ denotes ordinary ordinal addition. The symbol $\oplus$
denotes Hessenberg, or natural, sum: write the two Cantor normal forms using
a common decreasing list of exponents and add corresponding finite
coefficients. In particular,
\[
\alpha\oplus n=\alpha+n\quad(n<\omega),\qquad
\omega\oplus\omega=\omega\cdot2.
\]
Every infinite ordinal has a unique expression $\lambda+n$ with $\lambda$
a nonzero limit ordinal and $n<\omega$. Whenever a finite correction is
subtracted in this manuscript, it is subtracted \emph{inside that finite
coefficient}: $\lambda+(n-1)$, not an unspecified ordinal subtraction.

\subsection{Friendly sequences}

Two elements $x,y$ are \emph{friends} when $x\inc y$, that is, they are
incomparable. Following Vialard's definition, a bad sequence is
\emph{friendly}, or \emph{open-ended}, if at each choice the selected element
has a friend in the residual immediately before that choice
\cite[Definition~3.3]{Vialard2023}. A friend need not be selected later, and
one friend may witness several successive choices.

The corresponding tree rank satisfies
\begin{equation}\label{eq:f-rec}
\boxed{\displaystyle
\fot(P)=\sup_{\substack{x\in P\\\exists y\in P\ (x\inc y)}}
             \bigl(\fot(P_x)+1\bigr).}
\end{equation}
The empty supremum is zero. For a finite poset, this rank is exactly the
maximum length of a friendly sequence. For an infinite wpo, taking the
supremum of those finite lengths is generally \emph{not} the tree rank.

Let $G(P)$ be the simple incomparability graph on $P$. Define
\[
\str(P)=\{x\in P:\exists y\in P\ (x\inc y)\}.
\]
These are precisely the nonisolated vertices of $G(P)$.

\begin{lemma}\label{lem:strip}
For every wpo $P$,
\[
\fot(P)=\fot(\str(P))\leq\mot(\str(P))\leq\mot(P).
\]
The friendly rank is monotone under passage to an induced subposet.
\end{lemma}
\begin{proof}
A friendly sequence never selects a globally friendless element. Moreover,
every friend of one of its selected elements lies in $\str(P)$, since the
friendship relation is symmetric. Consequently the friendly trees of $P$
and $\str(P)$ have exactly the same sequences. The first inequality follows
by viewing this tree as a subtree of the bad-sequence tree on $\str(P)$.
For monotonicity, a friendly sequence in an induced subposet retains all its
witness friends when considered in the larger order. The same subtree
argument proves monotonicity of maximal order type.
\end{proof}

We will use the following established consequence of the published bound
for friendly rank. Its role is isolated so that the transfinite dependence
is explicit.

\begin{lemma}[Known limit-floor bound]\label{lem:floor}
If $\mot(\str(P))=\lambda+n$, where $\lambda$ is a nonzero limit ordinal and
$n<\omega$, then
\[
\lambda\leq\fot(P)\leq\lambda+n.
\]
In particular, $\mot(\str(P))=\lambda$ implies $\fot(P)=\lambda$.
\end{lemma}

This follows from the stronger bound
$\lambda+\lfloor n/2\rfloor\leq\fot(P)$ in
\cite[Theorem~4.5]{Vialard2023}, also stated in
\cite[Theorem~5.4.4]{Vialard2024}. We use neither the unrestricted disjoint-sum
identity in the former paper nor a deduction from that identity.

\section{The incomparability graph and ordinal sums}\label{sec:graph}

\begin{lemma}[Components are ordinally ordered]\label{lem:components}
The connected components of $G(P)$ are the factors in a canonical ordinal
sum decomposition
\[
P=\sum_{\xi<\kappa} C_\xi
\]
for an ordinal $\kappa$.
\end{lemma}
\begin{proof}
Elements in different components are comparable. If $a\inc b$ and $z$ lies
outside their component, the direction of comparison with $z$ must agree:
$a<z<b$ or $b<z<a$ would contradict $a\inc b$. Propagating this observation
along finite incomparability paths shows that every comparison between two
distinct components has the same direction. The resulting component order
is linear. An infinite descending sequence of components would yield an
infinite descending sequence in $P$ by choosing representatives, contrary
to the wpo property. Hence the component order is a well-order.
\end{proof}

\begin{theorem}[Transfinite ordinal-sum additivity]\label{thm:sum}
For any ordinal-indexed family of wpos $(A_\xi)_{\xi<\kappa}$,
\begin{equation}\label{eq:sum}
\fot\left(\sum_{\xi<\kappa}A_\xi\right)
   =\sum_{\xi<\kappa}\fot(A_\xi).
\end{equation}
Empty factors are permitted. The sum on the right is ordinary ordinal sum.
\end{theorem}
\begin{proof}
The ordinal sum is a wpo: a bad sequence cannot move to a strictly higher
factor, and an infinite nonincreasing sequence of ordinal indices is
eventually constant. An infinite bad sequence would therefore produce one
inside a factor.

Prove the formula by induction on the maximal order type of the entire
sum $P$. A selected $x\in A_\eta$ has a friend in $P$ exactly when it has a
friend in $A_\eta$. Its residual is
\[
P_x=\left(\sum_{\xi<\eta}A_\xi\right)+(A_\eta)_x.
\]
By \eqref{eq:mot-rec}, $\mot(P_x)<\mot(P)$, so the induction hypothesis
applies to this residual sum. Write
$\theta_\eta=\sum_{\xi<\eta}\fot(A_\xi)$. Then \eqref{eq:f-rec} gives
\[
\fot(P)=
\sup_{\substack{\eta<\kappa\\x\text{ has a friend in }A_\eta}}
\left(\theta_\eta+\fot((A_\eta)_x)+1\right).
\]
For a factor of positive friendly rank, continuity of addition in its
right argument makes the supremum over $x$ equal to
$\theta_\eta+\fot(A_\eta)$. A zero-rank factor contributes no terms. Taking
the supremum over the positive-rank factors gives precisely the ordinal sum
on the right of \eqref{eq:sum}; intervening or terminal zero terms do not
change it. If every factor has rank zero, both sides are zero.
\end{proof}

The binary case is already stated in the literature
\cite[Proposition~4.1(1)]{Vialard2023}. The proof above also treats limit
indices and zero-rank components, which will matter in applications.

\subsection{A maximal-vertex lemma}

A cut vertex of a connected graph is a vertex whose deletion leaves at
least two connected components. An isolated singleton is not a cut vertex.
The next elementary fact drives both the finite and infinite arguments.

\begin{lemma}[At most one cut maximum]\label{lem:cut}
In an incomparability-connected poset, at most one maximal element is a cut
vertex of the incomparability graph.
\end{lemma}
\begin{proof}
Suppose distinct maximal elements $u,v$ were both cut vertices. Choose a
component $D$ of $G(P)-u$ not containing $v$. Connectivity of the original
graph gives an $a\in D$ adjacent to $u$. There is no edge from $a$ to $v$,
so maximality of $v$ gives $a<v$.

Similarly choose a component $E$ of $G(P)-v$ not containing $u$, with a
vertex $b\in E$ adjacent to $v$. Then $b<u$. Since $b\inc v$, the vertex
$b$ cannot belong to $D$: otherwise its edge to $v$ would join $D$ to $v$
in $G(P)-u$. Thus $a$ and $b$ are in different components of $G(P)-u$, and
must be comparable. If $a<b$, then $a<b<u$, contrary to $a\inc u$. If
$b<a$, then $b<a<v$, contrary to $b\inc v$. Both alternatives are impossible.
\end{proof}

\begin{corollary}\label{cor:finite-noncut}
A finite incomparability-connected poset with at least two elements has a
maximal element whose deletion leaves the incomparability graph connected.
\end{corollary}
\begin{proof}
A finite poset with a unique maximal element has a greatest element: every
element lies below some maximal element. A greatest element is isolated in
the incomparability graph. The present poset therefore has at least two
maximal elements, and \cref{lem:cut} applies.
\end{proof}

\section{An exact formula for every finite poset}\label{sec:finite}

Write $c(P)$ for the number of components of $G(P)$ and set $c(\varnothing)=0$.

\begin{theorem}[Finite component formula]\label{thm:finite}
For every finite poset $P$,
\begin{equation}\label{eq:finite}
\boxed{\fot(P)=|P|-c(P).}
\end{equation}
There exists a friendly sequence attaining this length.
\end{theorem}
\begin{proof}
First let $C$ be a connected nonempty component with $r$ vertices. Any
friendly sequence selects distinct elements. Its final selected element
has an incomparable friend, and that friend survives the final deletion.
Thus at least one element of $C$ is left unselected, so its length is at
most $r-1$.

For the reverse inequality, if $r>1$, use
\cref{cor:finite-noncut} to select a maximal non-cut vertex $x$. It has a
friend because the graph is connected and nontrivial. Maximality makes
$C_x=C\setminus\{x\}$, still connected. Repeating this procedure leaves one
vertex after exactly $r-1$ friendly selections. The singleton case has rank
zero, as required.

Finally decompose $P$ into its finitely many ordinally ordered components.
By \cref{thm:sum},
\[
\fot(P)=\sum_C\fot(C)=\sum_C(|C|-1)=|P|-c(P).
\]
Witness sequences can be concatenated in descending component order. Once
we move to a lower component, a deletion may also remove the unused
singletons in higher components, but it does not remove the friends still
needed in the current lower component.
\end{proof}

\begin{corollary}[Graphic interpretation]\label{cor:forest}
For finite $P$, $\fot(P)$ equals the number of edges in any spanning forest
of $G(P)$, equivalently the rank of its graphic matroid.
\end{corollary}
\begin{proof}
A spanning tree of a component with $r$ vertices has $r-1$ edges. Summing
over the components gives \eqref{eq:finite}.
\end{proof}

The finite invariant is consequently unchanged by reversing every order
comparison: $P$ and its dual have the same incomparability graph. Every
value $r\in\{0,\ldots,n-1\}$ occurs on an $n$-element nonempty poset, for
example on $\Gamma_{r+1}+[n-r-1]$, where $\Gamma_k$ is a $k$-element
antichain and $[k]$ a $k$-element chain.

\subsection{Algorithm and certificate}

If the order relation is supplied as a reflexive transitive $n\times n$
matrix, form an incomparability edge exactly when neither direction of
comparison holds. A graph traversal counts its components in $O(n^2)$
time. The value of the rank then follows from a single subtraction. A Hasse
diagram input may first be transitively closed; the elementary dense method
costs $O(n^3)$.

A witness is constructed component by component, from highest to lowest:
\begin{lstlisting}
for C in incomparability_components_in_descending_order(P):
    while C contains more than one vertex:
        choose a maximal x in C for which G(C - {x}) is connected
        append x to the sequence
        C = C - {x}
\end{lstlisting}

At most two maximal candidates need to be tested at each stage, by
\cref{lem:cut}. Recomputing maxima and connectivity from a dense matrix
therefore gives an $O(n^3)$ witness algorithm. A separate checker verifies
the certificate by maintaining the actual residual of the original poset,
checking a surviving friend, and deleting the selected element's upset.
The archive implements this checker independently of the formula.

\begin{example}[Three scalar invariants already fail in the finite case]
The posets $\Gamma_2+\Gamma_2$ and $[2]\ds[2]$ both have four elements,
height two, and width two. Their incomparability graphs have respectively
two components and one component. Their friendly ranks are therefore two
and three. The stronger failure involving all four invariants in
\cref{thm:nonfunctional} requires infinite examples.
\end{example}

\section{The canonical finite upper tail}\label{sec:tail}

\begin{definition}\label{def:tail}
For a wpo $P$, define its \emph{finite upper tail} and \emph{core} by
\[
\tail{P}=\{x\in P:|\uparrow_Px|<\omega\},\qquad
\core{P}=P\setminus\tail{P}.
\]
The word ``tail'' does not assert that every core element lies below every
tail element. The set $T(P)$ is an upset, but generally not an ordinal-sum
factor of $P$.
\end{definition}

\begin{lemma}[Finiteness and stability]\label{lem:tailfinite}
The set $T(P)$ is a finite upset. For every finite subset $F\subseteq P$,
\[
T(P\setminus F)=T(P)\setminus F.
\]
The core has no maximal elements.
\end{lemma}
\begin{proof}
If $x\leq y$ and $\uparrow x$ is finite, then
$\uparrow y\subseteq\uparrow x$ is finite, proving upward closure.
Suppose $T(P)$ were infinite. An infinite subset of a wpo contains an
infinite strictly increasing sequence of distinct elements. To see this
last fact, enumerate a countable subset and apply the infinite Ramsey
theorem to the three possibilities for each pair of indices: increasing,
decreasing, or incomparable. The last two homogeneous possibilities are
excluded by the wpo property. An increasing sequence in $T(P)$ contradicts
the finiteness of the upset of its first element. Thus $T(P)$ is finite.

Removing finitely many points from an infinite upset does not make it
finite, and removing them from a finite upset leaves it finite. This proves
the stability identity. Finally, each $x\in\core{P}$ has infinitely many
successors in $P$; all but finitely many remain in the core. In particular
it has a strict successor there.
\end{proof}

\begin{lemma}[Removing a maximal element]\label{lem:maxdelete}
If $x$ is a maximal element of a wpo $P$, then
\[
\mot(P)=\mot(P\setminus\{x\})+1.
\]
\end{lemma}
\begin{proof}
Put $Q=P\setminus\{x\}$. A maximizing linear extension of $Q$, followed
by $x$, is a linear extension of $P$, giving the lower bound. Forgetting
all comparisons between $x$ and $Q$ gives the disjoint sum $Q\ds1$.
Every bad sequence in $P$ is bad in that weaker order. Hence
\[
\mot(P)\leq\mot(Q\ds1)=\mot(Q)\oplus1=\mot(Q)+1.
\]
\end{proof}

\begin{proposition}[The finite coefficient is intrinsic]\label{prop:tailtype}
If $P$ is infinite and $n=|T(P)|$, then
\begin{equation}\label{eq:tailtype}
\mot(P)=\lambda+n,\qquad\lambda=\mot(\core{P}),
\end{equation}
where $\lambda$ is a nonzero limit ordinal. Thus $n$ is exactly the finite
coefficient at the end of the Cantor normal form of $\mot(P)$.
\end{proposition}
\begin{proof}
The core is infinite because only finitely many points have been removed.
It has no maximal element. A linear extension of a poset with no maximal
elements has no last element, so a maximizing linear extension of the core
has nonzero limit type $\lambda$.

If the tail is nonempty, choose a maximal element in that finite upset. It
is maximal in $P$ because the tail is upward closed. Delete it and continue.
After $n$ such deletions the residual is exactly the core. Repeated use of
\cref{lem:maxdelete} gives \eqref{eq:tailtype}.
\end{proof}

This identifies the finite coefficient without selecting any particular
maximal linear extension. For instance, the core of $\omega\ds1$ is the
chain $\omega$, whereas the core of $(\omega\ds\omega)\ds1$ consists of two
incomparable infinite chains.

\subsection{When full successor rank is possible}

\begin{lemma}[Full rank forces maximal moves]\label{lem:fullmax}
Let $\mot(P)=\lambda+n$, with $\lambda$ a nonzero limit ordinal and
$0<n<\omega$. If $\fot(P)=\lambda+n$, there is an eligible maximal element
$x$ such that
\[
\fot(P\setminus\{x\})=\mot(P\setminus\{x\})=\lambda+(n-1).
\]
Consequently full rank implies $\fot(\core{P})=\lambda$.
\end{lemma}
\begin{proof}
A supremum equal to a successor ordinal is attained by one of its terms.
Equation~\eqref{eq:f-rec} therefore gives an eligible $x$ with
$\fot(P_x)=\lambda+(n-1)$. Since
$\fot(P_x)\leq\mot(P_x)<\mot(P)$, the maximal order type of $P_x$ must have
the same value.

Suppose $x<y$. No point of $P_x$ lies above $x$ or above $y$. A maximizing
linear extension of $P_x$, followed by $x$ and then $y$, is therefore a
linear extension of the induced subposet $P_x\cup\{x,y\}$. It follows that
\[
\mot(P)\geq\mot(P_x)+2=\lambda+(n+1),
\]
a contradiction. Thus $x$ is maximal, so $P_x=P\setminus\{x\}$.

Repeat this argument $n$ times. Every deleted point has a finite original
upset: at its deletion all its original strict successors have already
been deleted. The deleted set is therefore contained in $T(P)$ and has
its full cardinality $n$. The remaining order is the core and has friendly
rank $\lambda$.
\end{proof}

\section{The successor-defect theorem}\label{sec:connected}

The finite result loses one unit per connected component. In the infinite
case that unit may disappear through ordinal limit behavior. The theorem
below says precisely when it disappears.

\begin{lemma}[The exceptional cut maximum]\label{lem:unique-cut}
Let $C$ be an infinite incomparability-connected wpo with a unique maximal
element $u$. If $u$ is a cut vertex, then $C\setminus\{u\}$ has no maximal
elements, and $\mot(C)=\lambda+1$ for a nonzero limit $\lambda$.
\end{lemma}
\begin{proof}
Every component of $G(C)-u$ has a vertex adjacent to $u$, since the original
graph is connected. Suppose $z$ were maximal in $C\setminus\{u\}$. If
$z\inc u$, then $z$ would be a second maximal element of $C$. Hence $z<u$.
The component containing $z$ must be the last component of
$C\setminus\{u\}$ under \cref{lem:components}, since otherwise any point of
a later component would be above $z$.

There are at least two components after deleting the cut vertex. Choose
an earlier component $D$. Every $d\in D$ satisfies $d<z<u$, so no vertex
of $D$ is adjacent to $u$, a contradiction. Thus the remainder has no
maximal element. Its maximal order type is a nonzero limit, as in
\cref{prop:tailtype}, and \cref{lem:maxdelete} finishes the proof.
\end{proof}

\begin{lemma}[Deleting the tail without breaking connectivity]\label{lem:peel}
Let $C$ be infinite and incomparability-connected, with
$\mot(C)=\lambda+n$ and $1\leq n<\omega$. If $n\geq2$, a maximal vertex can be deleted while
preserving incomparability connectivity. Consequently one can delete
$n-1$ tail elements through friendly moves, leaving a connected order
$D$ with $\mot(D)=\lambda+1$. The final tail point of $D$ is also eligible,
and deleting it leaves $\core{C}$.
\end{lemma}
\begin{proof}
There is a maximal element because $n>0$ and $T(C)$ is nonempty. If there
are two or more maximal elements, \cref{lem:cut} supplies a non-cut one.
If there is only one, it cannot be a cut vertex when $n\geq2$, by
\cref{lem:unique-cut} and uniqueness of the finite coefficient. Each chosen
vertex has a friend because the current graph is infinite and connected.
Deleting it removes only that vertex and reduces the finite coefficient by
one. Tail stability keeps the core unchanged.

After $n-1$ deletions the remaining connected order has a tail consisting
of one point. That point is maximal and still has a friend. Its deletion
leaves the original core. No assertion that the core remains connected is
needed or intended.
\end{proof}

\begin{definition}[Core defect bit]\label{def:epsilon}
For an infinite wpo $P$, let $\lambda=\mot(\core{P})$ and define
\[
\eps(P)=
\begin{cases}
0,&\mot(\str(\core{P}))=\lambda,\\
1,&\mot(\str(\core{P}))<\lambda.
\end{cases}
\]
The stripping operation is performed \emph{after} the finite tail is removed.
\end{definition}

\begin{theorem}[Connected successor-defect classification]\label{thm:connected}
Let $C$ be a nonempty incomparability-connected wpo.
If $C$ is finite, then $\fot(C)=|C|-1$.
If $C$ is infinite, put $\lambda=\mot(\core{C})$ and $n=|T(C)|$. Then
\begin{equation}\label{eq:connected}
\boxed{\fot(C)=\lambda+\bigl(n-\eps(C)\bigr).}
\end{equation}
Here $\eps(C)\leq n$. In particular, if $n=0$ then
$\fot(C)=\mot(C)=\lambda$. If $n>0$, the only possible values are
$\lambda+(n-1)$ and $\lambda+n$.
\end{theorem}
\begin{proof}
The finite case is \cref{thm:finite}. If $n=0$, then
$\core{C}=C=\str(C)$, since an infinite connected graph has no isolated
vertex. The limit-floor bound gives $\fot(C)=\lambda$ and $\eps(C)=0$.

Suppose $n\geq1$. Apply \cref{lem:peel}, making $n-1$ friendly moves to a
connected order $D$ of maximal order type $\lambda+1$. Since
$\str(D)=D$, \cref{lem:floor} gives $\fot(D)\geq\lambda$. Each previously
made move adds one on the right to this lower bound. Hence
\begin{equation}\label{eq:two-values}
\lambda+(n-1)\leq\fot(C)\leq\lambda+n.
\end{equation}
There is no ordinal strictly between these two bounds.

If $\eps(C)=0$, then $\mot(\str(\core{C}))=\lambda$, and
\cref{lem:floor} gives $\fot(\core{C})=\lambda$. All $n$ deletions of the
tail described in \cref{lem:peel} are friendly, so
$\fot(C)\geq\lambda+n$, establishing equality with the upper bound.

Conversely, equality with that upper bound forces
$\fot(\core{C})=\lambda$ by \cref{lem:fullmax}. By \cref{lem:strip},
\[
\lambda=\fot(\core{C})
\leq\mot(\str(\core{C}))\leq\mot(\core{C})=\lambda.
\]
Thus full rank is possible only when $\eps(C)=0$. If $\eps(C)=1$, the
lower value in \eqref{eq:two-values} must occur.
\end{proof}

\begin{remark}[Why connectedness cannot be dropped]
An infinite chain has friendly rank zero and maximal order type equal to
its own order type. Its defect can therefore be arbitrarily large. The
connected theorem applies to each incomparability component, not to an
arbitrary order without first decomposing it.
\end{remark}

\begin{corollary}[A structural formula for all wpos]\label{cor:all}
Let $P=\sum_{\xi<\kappa}C_\xi$ be its incomparability-component decomposition.
For finite $C_\xi$, set $d_\xi=|C_\xi|-1$. For infinite $C_\xi$, set
\[
d_\xi=\mot(\core{C_\xi})+
          \bigl(|T(C_\xi)|-\eps(C_\xi)\bigr).
\]
Then $\fot(P)=\sum_{\xi<\kappa}d_\xi$.
\end{corollary}
\begin{proof}
Combine \cref{thm:sum,thm:connected}.
\end{proof}

The formula is intrinsic, but it is not a uniform decision procedure for
arbitrary presentations of infinite orders. It still requires maximal order
types of specified suborders and an equality test on the stripped core.

\section{Exact disjoint-sum laws}\label{sec:disjoint}

The incomparability graph of $A\ds B$, for nonempty $A,B$, is connected:
each vertex on one side is adjacent to every vertex on the other side.
Its tail and core satisfy
\begin{equation}\label{eq:sumcore}
T(A\ds B)=T(A)\ds T(B),\qquad
\core{(A\ds B)}=\core{A}\ds\core{B}.
\end{equation}
These identities follow immediately from the equality of the principal
upsets in a summand and in the disjoint sum.

\begin{theorem}[Disjoint sums, all cases]\label{thm:disjoint}
Let $A,B$ be nonempty wpos.
\begin{enumerate}[label=(\roman*)]
\item If both are finite, then
$\fot(A\ds B)=|A|+|B|-1$.
\item If both are infinite, then
\begin{equation}\label{eq:two-infinite}
\boxed{\fot(A\ds B)=\mot(A)\oplus\mot(B)=\mot(A\ds B).}
\end{equation}
There is no hypothesis that either maximal order type is a limit ordinal.
\item If $A$ is infinite and $B$ is finite of size $m>0$, write
$\mot(A)=\lambda+n$, with $\lambda=\mot(\core{A})$. Then
\begin{equation}\label{eq:mixed}
\boxed{\fot(A\ds B)=\lambda+\bigl(n+m-\eps(A)\bigr).}
\end{equation}
The symmetric rule applies when $A$ is finite and $B$ infinite.
\end{enumerate}
If one summand is empty, the result is simply the friendly rank of the
other summand.
\end{theorem}
\begin{proof}
The finite case follows from connectedness and \cref{thm:finite}.

If both summands are infinite, both cores in \eqref{eq:sumcore} are nonempty.
Every core element has a friend in the other core, so the stripped core is
the entire core. Its defect bit is zero. Apply \cref{thm:connected} and the
standard maximal-order-type law \eqref{eq:mot-rec}.

In the mixed case, the core of the disjoint sum is precisely $\core{A}$,
its tail has $n+m$ points, and its core defect bit is $\eps(A)$. Again apply
\cref{thm:connected}.
\end{proof}

The mixed formula depends on the \emph{cardinality} of the finite summand,
not on its internal comparisons. This is stronger than merely giving bounds
between consecutive ordinals. It also explains why information about the
infinite core, rather than just the total rank, is necessary.

\begin{corollary}[Finitely many nonempty summands]\label{cor:manysums}
For a finite disjoint family with at least two infinite summands, the
friendly rank is the natural sum of the maximal order types of all the
summands. With exactly one infinite summand $A$ and a positive total of
$m$ elements in the finite summands, formula~\eqref{eq:mixed} applies.
With all summands finite and at least two summands, the rank is one less
than their total cardinality.
\end{corollary}
\begin{proof}
With at least two infinite summands, their nonempty cores supply friends
to each other, even after all finite tails and finite summands are removed.
The core is therefore fully stripped. The other cases are applications of
\cref{thm:disjoint} after grouping the finite summands.
\end{proof}

\section{Explicit counterexamples and independent calculations}\label{sec:counterexamples}

The results in this section can be checked from the residual recurrence
and finite arguments without using the limit-floor bound. This gives a
separate check on the main examples.

\subsection{An ordinal and finitely many incomparable points}

\begin{proposition}\label{prop:ordinal-antichain}
For every ordinal $\alpha$ and every positive integer $m$,
\begin{equation}\label{eq:ordinal-antichain}
\fot(\alpha\ds\Gamma_m)=\alpha+(m-1).
\end{equation}
In particular $\fot(\alpha\ds1)=\alpha$.
\end{proposition}
\begin{proof}
For $\alpha=0$ this is the finite antichain formula. For $m=1$ and
$\alpha>0$, selecting $\beta<\alpha$ leaves $\beta\ds1$; selecting the
isolated point leaves the chain $\alpha$, whose friendly rank is zero.
Transfinite induction gives
\[
\fot(\alpha\ds1)=\sup\left(\{\beta+1:\beta<\alpha\}\cup\{1\}\right)
=\alpha.
\]
For $m\geq2$ and $\alpha>0$, induct on $\alpha$, with an inner induction on
$m$. A move in the antichain leaves $\alpha\ds\Gamma_{m-1}$ and contributes
$\alpha+(m-1)$ to the rank recurrence. A move at $\beta<\alpha$ in the
chain contributes $\beta+m$. The latter values are at most
$\alpha+(m-1)$: for successor $\alpha=\gamma+1$ their maximum is
$\gamma+m$, and for limit $\alpha$ their supremum is $\alpha$. The antichain
move attains the required bound.
\end{proof}

\subsection{A counterexample to the earlier printed identity}\label{sec:printed}

Proposition~4.1(2) in the published 2023 paper states
\begin{equation}\label{eq:printed}
\fot(A\ds B)=1+(\mot(A)-1)\oplus(\mot(B)-1).
\end{equation}
Take $A=\omega\ds\omega$ and $B=1$. We compute the actual left-hand side
explicitly.

For $k\geq1$, let $g_k=\fot(\omega\ds[k])$, and set $g_0=0$. A move in the
infinite chain leaves a finite chain disjoint from $[k]$. The corresponding
rank contributions have supremum $\omega$, by the finite formula. A move
at position $j$ in the finite chain leaves $\omega\ds[j]$. Thus
\[
g_k=\max\left(\omega,\max_{0\leq j<k}(g_j+1)\right)
=\omega+(k-1).
\]
Selecting position $k$ in either chain of $\omega\ds\omega$ leaves the
other infinite chain disjoint from $[k]$. Hence
\begin{equation}\label{eq:two-chains-direct}
\fot(\omega\ds\omega)=\sup_{k<\omega}(g_k+1)=\omega\cdot2.
\end{equation}

Now select the third, incomparable singleton. This is a friendly maximal
move, and the residual is $\omega\ds\omega$. Therefore
\[
\fot\bigl((\omega\ds\omega)\ds1\bigr)\geq\omega\cdot2+1.
\]
The maximal order type of that order is $\omega\cdot2+1$, so equality holds.
On the other hand, the right-hand side of \eqref{eq:printed} is
\[
1+(\omega\cdot2-1)\oplus(1-1)=1+\omega\cdot2=\omega\cdot2,
\]
using the infinite-ordinal subtraction convention also used in the source.
The limit ordinal $\omega\cdot2$ makes the discrepancy independent of any
choice about how to decrement a successor ordinal's finite coefficient.

\begin{proposition}\label{prop:printed-counterexample}
The unrestricted identity \eqref{eq:printed} fails for
$A=\omega\ds\omega$ and $B=1$. The actual value is $\omega\cdot2+1$, not
$\omega\cdot2$.
\end{proposition}

The 2024 thesis does not repeat \eqref{eq:printed}; its stated disjoint-sum
condition concerns two limit maximal order types. Accordingly this is an
explicit correction to the printed formula, not a claim about whether the
author had already recognized its limitation.

\subsection{All four invariant values still do not suffice}

Let
\[
L=\sum_{k<\omega}\Gamma_2,\qquad
A=L+\Gamma_2,\qquad B=(\omega+1)\ds1.
\]
The order $L$ consists of an $\omega$-sequence of two-point antichain levels;
every point at one level lies below every point at every later level.

The ordinary ordinal height $h(P)$ is the rank of the decreasing-sequence
tree, and ordinal width $w(P)$ is the rank of the antichain-sequence tree.
In the present examples every antichain has at most two elements and one
has two, so the ordinal width is exactly two. A maximal chain in $L$ has
type $\omega$, and a maximal chain in $A$ or $B$ has type $\omega+1$.
By \cref{thm:sum,thm:finite},
\[
\mot(L)=2\cdot\omega=\omega,\qquad
\fot(L)=1\cdot\omega=\omega.
\]
The following table records the resulting values; the two rightmost columns
also locate the information that is lost by those four invariants.

\begin{center}
\renewcommand{\arraystretch}{1.18}
\begin{tabular}{@{}lcccccc@{}}
\toprule
Order & $\mot$ & $h$ & $w$ & $\fot$ & $|T|$ & $\eps$\\
\midrule
$A=L+\Gamma_2$ & $\omega+2$ & $\omega+1$ & $2$ & $\omega+1$ & $2$ & $0$\\
$B=(\omega+1)\ds1$ & $\omega+2$ & $\omega+1$ & $2$ & $\omega+1$ & $2$ & $1$\\
\bottomrule
\end{tabular}
\end{center}

The formula for $\fot(B)$ is also an immediate instance of
\cref{prop:ordinal-antichain}. Both orders have exactly two maximal elements,
so adjoining the number of maximal elements to the list would not separate
them either.

\begin{theorem}[Failure of four-invariant compositionality]\label{thm:nonfunctional}
The friendly order type of $P\ds Q$ is not determined by the values of
$\mot,h,w,\fot$ on the two summands. Specifically,
\begin{equation}\label{eq:nonfunctional}
\fot(A\ds1)=\omega+3,\qquad
\fot(B\ds1)=\omega+2,
\end{equation}
although $A$ and $B$ agree in all four invariants.
\end{theorem}
\begin{proof}
For $A\ds1$, delete the two points of its top $\Gamma_2$ one at a time,
then delete the new incomparable singleton. All three moves are friendly:
the singleton witnesses the first two, and a point of $L$ witnesses the
third. Only the chosen point is removed at each step. The residual is $L$,
whose friendly rank is $\omega$. Thus $\fot(A\ds1)\geq\omega+3$; its
maximal order type is $\omega+3$, proving equality.

On the other side, $B\ds1=(\omega+1)\ds\Gamma_2$. Formula
\eqref{eq:ordinal-antichain} gives $\fot(B\ds1)=(\omega+1)+1=\omega+2$.
The input invariant values coincide, including those of the common
singleton summand, while the outputs differ. No rule depending only on
those input values can therefore compute the output in every case.
\end{proof}

The obstruction is not an exotic large ordinal. It occurs at the first
infinite ordinal, with width bounded by two and only two maximal elements.
In $A$, deleting the finite tail leaves an infinite supply of internal
friendships. In $B$, deleting the finite tail leaves a linear order with
none. The core defect bit records exactly this distinction.

\section{Consequences and worked families}\label{sec:applications}

\subsection{Finite multisets}

For completeness, the multiset ordering $\Mr(P)$ consists of finite
multisets over $P$, with
\[
m\leq_r n\quad\Longleftrightarrow\quad
m=n\ \text{or}\ 
\forall x\in m\setminus(m\cap n)\ 
\exists y\in n\setminus(m\cap n)\ (x<y).
\]
Intersection and subtraction respect multiplicities. This is the
Dershowitz--Manna-style multiset ordering, not the injective multiset
embedding order. Vialard's theorem gives
\begin{equation}\label{eq:multiset}
w(\Mr(P))=\omega^{\fot(P)}
\end{equation}
\cite[Theorem~3.4]{Vialard2023}. Applying the formulas proved above yields
\[
\boxed{w(\Mr(P))=\omega^{|P|-c(P)}\quad\text{for finite }P.}
\]
Even when $P$ is empty, this is correct: its multiset order contains just
the empty multiset and has width one. For two infinite wpos,
\[
w(\Mr(A\ds B))=\omega^{\mot(A)\oplus\mot(B)}.
\]
The four-invariant example also yields
$w(\Mr(A\ds1))=\omega^{\omega+3}$ but
$w(\Mr(B\ds1))=\omega^{\omega+2}$.

\subsection{Finite rectangular grids}

\begin{proposition}\label{prop:grids}
For $m,n\geq2$, with the product order on $[m]\times[n]$,
\[
\fot([m]\times[n])=mn-3,
\qquad
w(\Mr([m]\times[n]))=\omega^{mn-3}.
\]
\end{proposition}
\begin{proof}
The bottom corner $(0,0)$ and top corner $(m-1,n-1)$ are isolated in the
incomparability graph. Put $a=(0,n-1)$ and $b=(m-1,0)$; they are adjacent.
Every other noncorner vertex $(i,j)$ is adjacent to $a$ when
$i>0$ and $j<n-1$, or to $b$ when $i<m-1$ and $j>0$. The union of these two
conditions excludes only the bottom and top corners. Thus all noncorner
vertices form one component, giving exactly three components in total.
Apply \cref{thm:finite} and \eqref{eq:multiset}.
\end{proof}

If either factor is a singleton, the product is a chain and the rank is
zero; the displayed grid formula must not be used in that case.

\subsection{Ordinal sums of finite blocks}

For finite nonempty blocks $F_\xi$,
\[
\fot\left(\sum_{\xi<\kappa}F_\xi\right)
=\sum_{\xi<\kappa}\bigl(|F_\xi|-c(F_\xi)\bigr).
\]
This realizes genuinely transfinite ranks from finite graph calculations.
For instance,
\[
\fot\left(\sum_{k<\omega}\Gamma_2\right)=\omega,
\qquad
\fot\left(\sum_{\xi<\omega^2}\Gamma_3\right)
=2\cdot\omega^2=\omega^2.
\]
A block that is a chain contributes zero, regardless of its finite length.
The order of the blocks matters because ordinal addition is not
commutative. For example,
\[
\fot(L+\Gamma_3)=\omega+2,
\qquad
\fot(\Gamma_3+L)=2+\omega=\omega.
\]
These examples illustrate both finite losses and their absorption at a
later limit.

\subsection{Two successor behaviors}

Both $\omega\ds1$ and $(\omega\ds\omega)\ds1$ are connected and have one
point in the finite upper tail. In the first, the core is a chain, so the
friendly rank is the predecessor $\omega$ of its maximal order type
$\omega+1$. In the second, the core remains connected and fully stripped,
so the friendly rank retains the final successor:
\[
\fot((\omega\ds\omega)\ds1)=\omega\cdot2+1.
\]
This is precisely the distinction that a blanket ``subtract one'' rule
cannot capture.

\section{Computation and proof audit}\label{sec:audit}

\subsection{What was checked by execution}

The program \texttt{code/verify.py} enumerates naturally labeled posets on
$\{0,\ldots,n-1\}$: every strict comparison points from a smaller integer to
a larger integer. The generator adds a new largest label, choosing any
downset of the existing poset as its predecessor set. Every such poset is
generated exactly once. Every finite isomorphism class has at least one
natural labeling, but an isomorphism class can occur more than once.

For each generated poset, the program computes the rank independently by
memoizing the defining residual recurrence on vertex masks. It then checks
agreement with $n-c(P)$, constructs a witness by maximal non-cut deletions,
checks the witness against the original residual semantics, and tests the
at-most-one-cut-maximum assertion whenever the graph is connected.

\begin{center}
\begin{tabular}{@{}rrr@{}}
\toprule
Vertices & Naturally labeled posets checked & Connected graphs\\
\midrule
0 & 1 & 0\\
1 & 1 & 1\\
2 & 2 & 1\\
3 & 7 & 4\\
4 & 40 & 27\\
5 & 357 & 275\\
6 & 4,824 & 4,070\\
7 & 96,428 & 86,278\\
\midrule
Total & 101,660 & 90,656\\
\bottomrule
\end{tabular}
\end{center}

There were no failures. The verifier additionally checked all 25 grids with
$2\leq m,n\leq6$ using the residual-rank oracle. The machine-readable output
includes the full rank distributions, elapsed times, Python version, and
execution timestamp. The associated unit tests also exercise scrambled
vertex labels, invalid relations, empty orders, chains, antichains, ordinal
sums, disjoint sums, and invalid witness sequences.

\subsection{What finite checking cannot establish}

The exhaustive computations do not verify an infinite tree rank. For
example, finite suborders $[k]\ds[k]$ have friendly ranks $2k-1$, whose
supremum is only $\omega$, while \eqref{eq:two-chains-direct} gives
$\fot(\omega\ds\omega)=\omega\cdot2$. The additional successor in the
three-summand counterexample is even less visible from finite truncation.
The transfinite statements depend on the mathematical proofs, not on
extrapolation from the finite data.

\subsection{Dependency boundaries}

\begin{center}
\small
\begin{tabular}{@{}p{0.32\linewidth}p{0.62\linewidth}@{}}
\toprule
Result & Dependencies and verification status\\
\midrule
Finite component formula & Elementary component and cut-vertex lemmas;
rank recurrence; exhaustive finite check.\\[3pt]
Canonical finite tail & Standard wpo chain property, maximal linear
extensions, natural-sum law, and the supplied deletion proof.\\[3pt]
Connected classification & Supplied graph and tail arguments plus the
previously published limit-floor bound.\\[3pt]
General disjoint-sum laws & Connected classification and the explicit
core identities.\\[3pt]
Printed-identity counterexample & Direct recurrences for finite and
infinite chains; no use of the disputed identity or the limit-floor bound.\\[3pt]
Four-invariant counterexample & Finite formula, transfinite ordinal-sum
additivity, and direct ordinal-plus-antichain calculation.\\[3pt]
Multiset consequences & The supplied rank calculations and Vialard's
multiset-width theorem.\\
\bottomrule
\end{tabular}
\end{center}

The new proofs are ordinary mathematical arguments rather than formal
proof objects. In particular, a future formalization should preserve three
points that are easy to obscure: every friend belongs to the \emph{current}
residual; every full-rank successor step is forced to select a maximal
element; and the core need not remain connected after the last tail point
is deleted.

\section{What remains to investigate}\label{sec:future}

The results answer the finite problem completely and give an exact
structural classification for all incomparability components. They also
supply a general disjoint-sum law once the core defect bit is available.
The countable example proves that unrestricted four-invariant
compositionality is not a viable substitute for that extra information.

A concrete next question is whether a useful class of presentations admits
compositional calculation of
\[
P\longmapsto\bigl(\mot(P),\ |T(P)|,\
                  \mot(\str(\core{P}))\bigr).
\]
For example, one can ask this for orders built from ordinal constants,
finite posets, ordinal sums, disjoint sums, and Cartesian products. The
present paper settles the friendly-rank step once those data and the
component decomposition are known, but does not provide every transition
rule for that presentation language.

Another question is whether other common operations preserve a tractable
core-saturation predicate. The disjoint sum of two infinite orders forces
saturation automatically; an infinite order plus a finite disjoint summand
preserves the defect bit of the infinite core. Products and higher-order
constructions may require additional structural information. The finite
grid calculation is an exact result for one product family, not a claimed
classification of arbitrary products.

Finally, the elementary proofs of the two explicit counterexamples are
small enough to make good formalization targets. A verified library of
well-founded tree ranks and ordinal sums could certify them independently
before attempting the more general connected-component classification.

\appendix
\section{Reproducing the artifacts}\label{sec:reproduce}

The Python implementation uses only the standard library and is written
for Python 3.9 or later. From the archive's top directory, run:
\begin{lstlisting}
python code/verify.py --max-n 7 --grid-max 6 \
    --output results/verification.json
python -m unittest discover -s code -v
\end{lstlisting}
The reference rank implementation can be exponential in the number of
vertices; it is intentionally retained as an independent oracle. The
component formula is the efficient method. Increasing the exhaustive
bound increases the number of generated posets rapidly.

To rebuild the PDF, run \texttt{python build.py}, or use:
\begin{lstlisting}
latexmk -pdf -interaction=nonstopmode -halt-on-error \
    friendly_order_types.tex
\end{lstlisting}
The source uses the freely available \texttt{newpx} LaTeX packages for text
and mathematics. Font binaries and third-party papers are not included.
The archive also contains a source-status note and a short proof-audit
checklist. The supplied \texttt{.bib} file and \texttt{.bbl} file allow both
normal BibTeX builds and inspection of the rendered references.

\bibliographystyle{plain}
\bibliography{references}
\end{document}
